% Task 12 insertable theory supplement. Standalone source; no venue formatting.
% No TeX engine was available locally; PDF typesetting is not yet verified.
\documentclass[11pt]{article}
\usepackage{amsmath,amssymb,amsthm}
\usepackage[margin=1in]{geometry}
\usepackage[hidelinks]{hyperref}
\newtheorem{proposition}{Proposition}
\newtheorem{theorem}[proposition]{Theorem}
\newtheorem{remark}[proposition]{Remark}
\newcommand{\norm}[1]{\left\lVert#1\right\rVert}
\title{Conditional Error and Cost Analysis of Change-Gated Streaming Inference}
\author{SOKKANAEM working supplement: task 12}
\date{6 September 2026}
\begin{document}
\maketitle

\paragraph{Scope.}
This supplement analyzes the frozen v11 implementation without changing its
weights or evaluating the reserved test set. The bounds below are conditional
perturbation results, not a certificate of depth accuracy. Update/copy gating
has precedent in Skip RNN~\cite{campos2018}; the identity alone is not claimed
as a new mathematical contribution. We distinguish state preservation, cached
readout approximation, spatial-context omission, and ground-truth error.

\section{Implemented recurrence and zero step}
Vectorize the inner-channel and state coordinates of one temporal SSM.
For a shared exogenous token sequence, write
\begin{equation}\label{eq:recurrence}
 F_t(h)=\Phi_t h+q_t,\qquad
 \Phi_t=\operatorname{diag}(e^{-\lambda_j\Delta_{t,j}}),\quad
 q_{t,j}=\Delta_{t,j}(B_tx_t)_j,
\end{equation}
where $\lambda_j>0$ and $\Delta_{t,j}\ge0$. A channel's step size is repeated
over its state coordinates. The code uses $A=-\exp(A_{\log})$ and a softplus
step size. Norms are Euclidean and induced operator norms unless specified.

\begin{proposition}[Binary zero-step identity]\label{prop:identity}
For a diagonal binary mask $M_t$, multiplying the step size by the mask gives
\begin{equation}\label{eq:gate}
 \widetilde h_t=M_tF_t(\widetilde h_{t-1})+(I-M_t)\widetilde h_{t-1}.
\end{equation}
Thus each inactive hidden-state coordinate is unchanged in real arithmetic.
\end{proposition}
\begin{proof}
On an active coordinate the original update is used. On an inactive coordinate
the transition is $e^0=1$ and the input term is zero. These cases give the
stated expression. Fractional masks do not generally have this update/copy
interpretation because $e^{m\Delta A}\ne m e^{\Delta A}+1-m$.
\end{proof}

The input term is first order, not exact zero-order hold (ZOH). With the input
held constant, the ZOH coefficient is $(1-e^{-\lambda\Delta})/\lambda$;
see the ZOH definition in Mamba~\cite{gu2024}. Its difference from the
implemented coefficient obeys
\begin{equation}\label{eq:zoh}
 \left|\Delta b-\frac{1-e^{-\lambda\Delta}}{\lambda}b\right|
 \le \frac{\lambda\Delta^2}{2}|b|.
\end{equation}
Indeed, the coefficient difference is
$\int_0^\Delta(1-e^{-\lambda s})\,ds\le\int_0^\Delta\lambda s\,ds$.
This local discretization error is distinct from skipping error: both paths
below use the same implemented recurrence. State identity also does not imply
readout identity, since the current input affects $C,x,z,D$ and the residual.
Finite CPU/CUDA tests supplement the real-arithmetic statement; they do not
establish bitwise identity for exceptional values or all floating-point scan
reassociations.

\section{Skipping defects and state error}
\begin{theorem}[Shared-input defect propagation]\label{thm:defect}
Let $h_t=F_t(h_{t-1})$ and let $\widetilde h_t$ satisfy
\eqref{eq:gate}, using the same $\Phi_t,q_t$. Define
\begin{equation}
 e_t=\widetilde h_t-h_t,\qquad
 r_t=(I-M_t)((I-\Phi_t)\widetilde h_{t-1}-q_t),\qquad d_t=\norm{r_t}.
\end{equation}
Then $e_t=\Phi_t e_{t-1}+r_t$. If $\norm{\Phi_t}\le\rho_t$, then
\begin{equation}\label{eq:bound}
 \norm{e_T}\le\left(\prod_{j=1}^T\rho_j\right)\norm{e_0}
 +\sum_{i=1}^T\left(\prod_{j=i+1}^T\rho_j\right)d_i.
\end{equation}
Empty products equal one. In particular, if $\rho_t\le\rho<1$ and
$d_t\le\epsilon$, the bound is
$\rho^T\norm{e_0}+\epsilon(1-\rho^T)/(1-\rho)$.
\end{theorem}
\begin{proof}
Add and subtract $\Phi_t\widetilde h_{t-1}+q_t$ in the gated update and
subtract the dense update. Iteration gives a sum of propagated defects.
Submultiplicativity and the triangle inequality give \eqref{eq:bound};
the uniform case follows by summing a geometric series.
\end{proof}

A uniform lower bound $\lambda_j\ge\lambda_{\min}>0$ and
$\Delta_{t,j}\ge\Delta_{\min}>0$ would yield
$\rho=e^{-\lambda_{\min}\Delta_{\min}}<1$ for the ungated comparison
transition. The actual gated transition equals one on skipped coordinates.
Without a useful uniform contraction bound, the nonexpansive estimate
$\norm{e_T}\le\norm{e_0}+\sum_i d_i$ remains available.

\paragraph{Counterexample: zero input change.}
For $F(h)=h/2+1$, equal zero initial states and all updates skipped,
$\widetilde h_t=0$ while $h_t=2(1-2^{-t})$, despite constant inputs.
The error at $t=4$ is $1.875$. An all-active first frame does not remove
the subsequent discrepancy. Thus a small pixel-change threshold alone does
not bound the omitted innovation, even when the input change is exactly zero.

\paragraph{Different token trajectories.}
If the gated branch uses $\widetilde\Phi_t,\widetilde q_t$, define its skip
defect using those parameters. The error recurrence acquires the extra term
$(\widetilde\Phi_t-\Phi_t)\widetilde h_{t-1}+\widetilde q_t-q_t$.
Its norm is at most
$\norm{\widetilde\Phi_t-\Phi_t}\norm{\widetilde h_{t-1}}
+\norm{\widetilde q_t-q_t}$. This term cannot be dropped for deeper blocks
whose inputs depend on previously approximated features.

\begin{theorem}[Whole-network conditional perturbation]\label{thm:network}
Fix the frame sequence and the masks selected by the sparse trajectory.
Let $z$ comprise all temporal states and block output caches, with the dense
comparison overwriting an identically shaped set of caches at every step.
Let $F_t$ be the all-active network step and $G_t$ the sparse step. Put
$d_t^{\rm net}=\norm{G_t(\widetilde z_{t-1})-F_t(\widetilde z_{t-1})}$.
If $F_t$ is $L_t$-Lipschitz on a region containing both states, then
\begin{equation}
 E_t^{\rm net}\le L_tE_{t-1}^{\rm net}+d_t^{\rm net}.
\end{equation}
Equation~\eqref{eq:bound} holds with $L_t,d_t^{\rm net}$ in place of
$\rho_t,d_t$. For dense and sparse output maps $P_t,\widetilde P_t$, if
$P_t$ is $L_{P,t}$-Lipschitz and
$\eta_t=\norm{\widetilde P_t(\widetilde z_{t-1})-P_t(\widetilde z_{t-1})}$,
then $\norm{\widetilde y_t-y_t}\le L_{P,t}E_{t-1}^{\rm net}+\eta_t$.
\end{theorem}
\begin{proof}
Add and subtract $F_t(\widetilde z_{t-1})$, or respectively
$P_t(\widetilde z_{t-1})$, and apply the assumed Lipschitz inequality.
Iterate the state bound.
\end{proof}

No whole-network $L_t<1$ or practically small global constant is certified
here. Negative diagonal $A$ alone is insufficient: the elementary map
$F(h)=0.5h+u(h)$ with $u(h)=h$ expands by $1.5$. This is a counterexample
to an inference about feedback, not an identification of v11 with that map.
Computing a counterfactual dense defect also costs work; the detector does
not provide it as a free online certificate. Threshold-induced mask changes
under perturbed images are outside this fixed-mask comparison.

For $n$ fixed pixels with ground truth $d_i\ge d_*>0$, the reverse triangle
inequality gives the additional, no-fit statement
\begin{equation}
 |\operatorname{AbsRel}(\widetilde y,d)-\operatorname{AbsRel}(y,d)|
 \le \frac1n\sum_i\frac{|\widetilde y_i-y_i|}{d_i}
 \le \frac{\norm{\widetilde y-y}_1}{nd_*}.
\end{equation}
This does not remove the dense reference's own ground-truth error. Separately
refitted scale--shift alignment requires conditioning assumptions of its own;
threshold accuracy metrics require margin assumptions. No new depth cutoff
has been introduced into the evaluation protocol.

\section{Cached readout and refresh}
\begin{proposition}[Conditional first-block cache bound]\label{prop:cache}
Consider a pixel-gated patch with $C p^2$ pixel values that remains inactive
for $a=t-s$ steps after its last update. Suppose its embedding is $u=WI+b$
and the state-fixed block readout $R(u,h_s)$ is $L_R$-Lipschitz in $u$ on
the relevant region. Then
\begin{equation}\label{eq:cache}
 \norm{R(u_t,h_s)-R(u_s,h_s)}
 \le L_R\norm{W}\,a\sqrt{Cp^2\tau_{\rm on}}.
\end{equation}
\end{proposition}
\begin{proof}
Inactivity implies patch MSE at most the applicable hysteresis threshold,
which is at most $\tau_{\rm on}$. Dilation and forced activation only turn
zeros into ones and preserve this implication for a final zero mask.
Each consecutive patch difference is at most $\sqrt{Cp^2\tau_{\rm on}}$.
Sum these differences, apply the embedding norm and then the readout bound.
The cached readout at $s$ used the already updated state $h_s$, which remains
unchanged during the inactive interval.
\end{proof}

The readout includes normalization, residual and input-dependent projections.
This first-block, pixel-MSE statement does not directly apply to deeper-layer
tokens or GMC relative-feature scores. Spatial active-subsequence computation
has an additional context defect
$\norm{S_A(u_A)-[S(u)]_A}$, even on active tokens. For example, the spatial
recurrence $s_i=\rho s_{i-1}+b_i$ with $b_1=1,b_2=0,s_0=0$ produces $s_2=\rho$;
omitting token 1 gives zero, even if the current inputs equal past inputs.
Current GMC warps the comparison image, not the hidden-state or output caches.

\begin{theorem}[Refresh is not state reset]\label{thm:refresh}
An all-active keyframe makes $G_t=F_t$ at the same incoming extended state;
hence $d_t^{\rm net}=0$, but generally
$E_t^{\rm net}\le L_tE_{t-1}^{\rm net}\ne0$.
For the shared-input SSM of Theorem~\ref{thm:defect}, assume a common
$\rho<1$, skip defects at most $\epsilon$, and a refresh every $K$ steps.
If $B_n$ is the error immediately after refresh $n$, then
\begin{align}
 B_{n+1}&\le\rho^KB_n+\rho\epsilon\frac{1-\rho^{K-1}}{1-\rho},\\
 \limsup_{n\to\infty}B_n&\le
 \frac{\rho\epsilon(1-\rho^{K-1})}{(1-\rho)(1-\rho^K)}.
\end{align}
\end{theorem}
\begin{proof}
All-active branches recompute block outputs but pass the carried temporal
states to the recurrence. Thus only the same-state local defect vanishes.
For the scalar norm recurrence, propagate $K-1$ bounded defects and apply
the final zero-defect refresh. Each injected term receives at least one
factor $\rho$. Sum the resulting cycle recurrence as a geometric series.
\end{proof}

At $j\le K-1$ steps after a refresh, the bound is
$\rho^jB_n+\epsilon(1-\rho^j)/(1-\rho)$. Actual errors need not be monotone
in $K$: changing $K$ also changes the trajectories and defects. Keyframes
bound cache age by $K-1$, not dense-history error by zero. Resetting only one
path to zero is a different intervention. True state synchronization would
require the reference state or replay. Rolling refresh bounds per-patch age
but is not an all-active whole-network refresh.

\section{Fixed cost, refresh frequency and overhead}
\begin{proposition}[Conditional pipeline speedup]\label{prop:cost}
Assume an additive cost model with dense reference $C_d=F+V>0$ and mean
sparse cost $\overline C_s=F+\overline a_{\rm eff}V+\overline H$, where
$F,V,\overline H\ge0$ and $\overline C_s>0$.
Define $f=F/C_d$ and $\omega=\overline H/C_d$.
Then
\begin{equation}
 S=\frac{C_d}{\overline C_s}
 =\frac1{f+(1-f)\overline a_{\rm eff}+\omega},\qquad
 S>1\ \Longleftrightarrow\ \omega<(1-f)(1-\overline a_{\rm eff}).
\end{equation}
For $f>0$, $S\le1/f$.
\end{proposition}
\begin{proof}
Divide by $C_d$ and compare the positive denominator with one.
Dropping its other nonnegative terms yields the upper bound.
\end{proof}

This is an application of Amdahl-style fixed-cost reasoning~\cite{amdahl1967},
not a new universal speed law. Detector, gather/scatter, cache-copy and launch
overheads must be counted without double counting. Keyframe and fallback
events are a union: $a_{{\rm eff},t}=1$ on either event, and otherwise equals
the normal activity. Padding and nonlinear kernel costs must be modeled
separately if active fraction does not scale variable work linearly.

For $T$ frames starting at frame index zero, the keyframe fraction is
$\lceil T/K\rceil/T$, even on a static stream. At $T=256,K=30$ this is
$9/256$, not exactly $1/30$. In a long stream without fallback and with
constant nonkeyframe mean $a$, the effective activity tends to
$a+(1-a)/K$. If a verified cache-only coefficient
$c=L_R\norm{W}\sqrt{Cp^2\tau}>0$ and tolerance $E_c$ are available,
\eqref{eq:cache} gives the sufficient condition
$K\le1+\lfloor E_c/c\rfloor$. In the nondegenerate fixed-cost model,
a target speedup $S_*>1$ requires
\begin{equation}
 D=1/S_*-f-(1-f)a-\omega>0,\qquad
 K\ge\left\lceil\frac{(1-f)(1-a)}{D}\right\rceil.
\end{equation}
Nonoverlap means these sufficient error conditions cannot certify a feasible
$K$; it is not a proof that no accurate fast policy exists.

The existing K30 synchronized profile assigns about $69.4\%$ of its time
to image read/decode/preprocessing. Holding that component fixed and deleting
all remaining profiled time gives an idealized ceiling of about $1.44\times$
relative to that same profile. This is neither a measured speedup nor a bound
for a redesigned input path or GPU-resident pipeline. Main matched-subset
latencies remain $7.496$ ms sparse and $7.521$ ms dense.
Persistent temporal states are densely allocated:
$bBN\sum_{\ell\in\mathrm{temporal}}P_\ell S_\ell$ bytes, independent of activity.
For v11 this is $12$ MiB in FP32; output caches add storage.

\paragraph{Reproducibility and limitations.}
The accompanying \texttt{tests/test\_theory.py} checks algebraic identities,
counterexamples and actual SSM/cache/refresh behavior. Numerical checks do
not replace these proofs or certify global constants for the trained model.
The executable audit and existing measurements are indexed in
\texttt{paper/THEORY\_12.md}. Full-manuscript integration, further model
improvement, device validation and reserved-test inference are separate work.

\begin{thebibliography}{9}
\bibitem{gu2024} A. Gu and T. Dao.
Mamba: Linear-Time Sequence Modeling with Selective State Spaces.
arXiv:2312.00752v2, 2024.
\url{https://arxiv.org/html/2312.00752v2}.
\bibitem{campos2018} V. Campos, B. Jou, X. Gir\'o-i-Nieto, J. Torres and S.-F. Chang.
Skip RNN: Learning to Skip State Updates in Recurrent Neural Networks.
ICLR, 2018. \url{https://arxiv.org/html/1708.06834v3}.
\bibitem{amdahl1967} G. M. Amdahl.
Validity of the single processor approach to achieving large scale computing capabilities.
AFIPS Spring Joint Computer Conference, pp. 483--485, 1967.
\url{https://doi.org/10.1145/1465482.1465560}.
\end{thebibliography}
\end{document}
